% Part: first-order-logic
% Chapter: proof-systems
% Section: natural-deduction

\documentclass[../../../include/open-logic-section]{subfiles}

\begin{document}

\iftag{FOL}
      {\olfileid{fol}{prf}{ntd}}
      {\olfileid{pl}{prf}{ntd}}

\olsection{自然演绎}

自然演绎是一种旨在反映实际推理（尤其是数学家所使用的那种严密规范的推理）的推导系统。
实际推理遵循许多“自然”的模式。
例如，分情形证明允许我们在一个析取前提的基础上确立一个结论，方式是通过表明该结论从两个析取支中的任何一个都能推出。
间接证明则让我们可以通过证明结论的否定会导致矛盾来确立该结论。
条件证明通过证明后件能从前件推出，来确立一个条件句“如果……那么……”。
自然演绎便是其中一些自然推理的形式化。
每个逻辑联结词和量词都对应两条规则，一条引入规则和一条消去规则，它们各自对应一种这样的自然推理模式。
例如，$\Intro{\lif}$ 对应条件证明，而 $\Elim{\lor}$ 对应分情形证明。
特别简单的一条规则是 $\Elim{\land}$，它允许从 $!A \land !B$ 推出~$!A$（或 $!B$）。

自然演绎区别于其他推导系统的一个特点是它对假设的运用。
自然演绎中的一个推导是一个公式树。
单个公式位于公式树的根部，而树的“叶子”则是推导出结论所依据的公式。
在自然演绎中，有些叶公式在推导过程中扮演了角色，但在推导到达结论时就被“用完”了。
这对应于实际推理中引入仅暂时生效的假设的做法。
例如，在分情形证明中，我们假设每个析取支为真；在条件证明中，我们假设前件为真；在间接证明中，我们假设结论的否定为真。
这种引入假设并在确立中间步骤时将其消去的做法，是自然演绎的标志。
自然演绎推导中位于叶子的公式称为\emph{假设}，而某些推理规则可以“解除”它们。
例如，如果我们有一个从包含~$!A$ 在内的一些假设推出~$!B$ 的推导，那么 $\Intro{\lif}$ 规则就允许我们推断出~$!A \lif !B$，并解除任何形式为~$!A$ 的假设。
（为了追踪哪些假设在哪些推理步骤被解除，我们为推理步骤及被其解除的假设标上相同的数字。）
推导结束时仍未被解除的假设，共同足以保证结论为真，因此一个推导确立了其未解除的假设是其结论的语义后承。

基于自然演绎的关系 $\Gamma \Proves !A$ 成立，当且仅当存在一个推导，其中 $!A$~是树中最后一个语句，且每个未被解除的叶公式都属于~$\Gamma$。
$!A$~是自然演绎中的定理，当且仅当存在一个推导，其中 $!A$~是最后一个语句且所有假设均被解除。
例如，以下是一个表明 $\Proves (!A
\land !B) \lif !A$ 的推导：

\begin{prooftree}
  \AxiomC{$\Discharge{!A \land !B}{1}$}
  \RightLabel{\Elim{\land}}
  \UnaryInfC{$!A$}
  \DischargeRule{\Intro{\lif}}{1}
  \UnaryInfC{$(!A \land !B) \lif !A$}
\end{prooftree}

标签~$1$ 表明假设 $!A \land !B$ 在 \Intro{\lif} 推理步骤处被解除。

一个集合~$\Gamma$ 是不一致的，当且仅当在自然演绎中 $\Gamma \Proves \lfalse$。
规则 \FalseInt{} 使得从一个不一致的集合中可以推导出任何语句。

自然演绎系统由 Gerhard Gentzen（根岑）和 Stanisław Jaśkowski（亚斯科夫斯基）在 20 世纪 30 年代创立，后来由 Dag Prawitz（普拉维茨）和 Frederic Fitch（菲奇）加以发展。
由于其推理反映了自然的证明方法，它受到哲学家的青睐。
Fitch 发展的版本常在入门逻辑教科书中使用。
在逻辑哲学中，自然演绎的规则有时被认为给出了逻辑算子的意义（“证明论语义学”）。

\end{document}